
\subsubsection{2.1 Shortcut Learning in Deep Neural Networks}

Geirhos et al. (2020) provide a taxonomy of shortcut learning across vision, language, and medical imaging domains, characterizing shortcuts as decision rules that fail under distribution shift.

\subsubsection{2.2 Simplicity Bias and Feature Learning}

Shah et al. (2020) formalized simplicity bias: neural networks trained with gradient descent preferentially learn features that are linearly separable in the input space. Hermann and Lampinen (2020) extended this analysis to show that models learn complex features more slowly than simple ones.

\subsubsection{2.3 Gradient Starvation}

Pezeshki et al. (2021) identified gradient starvation as the mechanism by which dominant features suppress learning of minority features. When one feature explains most of the variance in the loss, gradients flowing to alternative feature pathways diminish.

\subsubsection{2.4 Group Robustness Methods}

Sagawa et al. (2020) introduced Group DRO and the Waterbirds benchmark. Liu et al. (2021) proposed Just Train Twice (JTT). Kirichenko et al. (2023) showed that last layer retraining (DFR) recovers most of the worst-group accuracy lost to shortcuts.

\subsubsection{2.5 Positioning This Work}

Prior work describes what features are learned (simplicity bias), why shortcuts persist (gradient starvation), and how to correct them (Group DRO, JTT, DFR). This work contributes the temporal dimension: when classifier commitment accelerates.

